% The chapter's introduction (from myChessNotebook, chapter 6).
Most chess books are not very useful for my own training.  
Many authors either present endless variations or write as if their conclusions were absolute truth.  
Even when some statements are doubtful (for example, ``queen and knight is always better than queen and bishop''), practical chess over the board is much more dynamic than final-looking analysis.

Real chess is played move by move.  
In this chapter, I will replay grandmaster games one move at a time and follow one side closely.  
During each playout, I will record what I am thinking: my plan, my opponent's plan, and the key practical decisions.  
Whenever possible, I will avoid long side-variations and focus on the main ideas.

The goal is to understand my thinking habits in different types of positions.  
Do I repeat the same thought patterns?  
Do I miss critical ideas or candidate moves?

Grandmaster games plus computer analysis will help me identify the blind spots in my decision process.

To keep the notes clear, I will use three distinct styles:
\begin{itemize}
  \item My own thoughts: \mythoughts{italic serif text}. 
  \item My reflection after I see the computer line: \myreflection{teal text}. 
\end{itemize}
